% L-BFGS vs 精确 Newton vs GD：收敛速度对比
% 横轴：迭代步数；纵轴：梯度范数 \|\nabla f\|（对数坐标）
% 三条曲线：GD（慢）/ L-BFGS（接近 Newton）/ Newton（最快）
\documentclass[tikz,border=4pt]{standalone}
\usepackage{ctex}
\usepackage{amsmath,amssymb}
\usepackage{pgfplots}
\pgfplotsset{compat=1.18}

\begin{document}
\begin{tikzpicture}

\begin{axis}[
  width=10cm, height=8cm,
  xlabel={迭代步数 $t$},
  ylabel={$\log_{10}\|\nabla f(\theta^t)\|$},
  xlabel style={font=\small},
  ylabel style={font=\small},
  tick label style={font=\footnotesize},
  xmin=0, xmax=40,
  ymin=-9, ymax=1,
  grid=both, grid style={gray!20},
  axis lines=left,
  legend style={
    font=\footnotesize,
    at={(0.98, 0.97)},
    anchor=north east,
    draw=gray!50,
    fill=white,
    rounded corners=2pt,
  },
  title={二阶方法 vs 一阶方法：收敛速度对比},
  title style={font=\small\bfseries},
]

% ─────────────────────────────────────────────────────
% 曲线 1：梯度下降（GD）—— 线性收敛，较慢
% 条件数 κ=100，线性收敛率 ≈ (κ-1)/(κ+1)≈0.98
% log10(0.98^t) = t*log10(0.98) ≈ -0.00877t
% ─────────────────────────────────────────────────────
\addplot[
  red, very thick,
  domain=0:40,
  samples=41,
] {-0.0877*x};
\addlegendentry{梯度下降（GD，线性收敛，$\kappa=100$）}

% ─────────────────────────────────────────────────────
% 曲线 2：L-BFGS —— 超线性收敛
% 前 5 步类似 GD，之后加速，最终接近 Newton 速率
% 用分段近似：先慢后快
% ─────────────────────────────────────────────────────
\addplot[
  blue, very thick,
  smooth,
] coordinates {
  (0,  0.0)
  (2,  -0.20)
  (4,  -0.55)
  (6,  -1.10)
  (8,  -1.80)
  (10, -2.70)
  (12, -3.70)
  (14, -5.00)
  (16, -6.50)
  (18, -8.00)
  (20, -8.80)
};
\addlegendentry{L-BFGS（超线性收敛，$m=10$ 历史）}

% ─────────────────────────────────────────────────────
% 曲线 3：精确 Newton —— 二次收敛（最快）
% 一旦进入收敛域，每步平方误差
% ─────────────────────────────────────────────────────
\addplot[
  green!50!black, very thick, dashed,
  smooth,
] coordinates {
  (0,  0.0)
  (2,  -0.30)
  (4,  -0.80)
  (6,  -2.0)
  (8,  -4.5)
  (10, -8.8)
};
\addlegendentry{精确牛顿法（二次收敛）}

% ─────────────────────────────────────────────────────
% 标注区域
% ─────────────────────────────────────────────────────
% L-BFGS 接近 Newton 的区域标注
\draw[<-, blue!70, thin]
  (axis cs:16, -6.5) -- (axis cs:22, -5.5)
  node[right, font=\tiny, blue!70, align=left]
    {L-BFGS 超线性加速\\（历史 $m$ 步 Hessian 近似）};

% GD 收敛慢标注
\draw[<-, red!70, thin]
  (axis cs:30, -2.63) -- (axis cs:24, -1.5)
  node[left, font=\tiny, red!70, align=right]
    {GD：$O\!\left((\tfrac{\kappa-1}{\kappa+1})^t\right)$\\受条件数 $\kappa$ 制约};

% Newton 二次收敛标注
\draw[<-, green!50!black, thin]
  (axis cs:9, -6.5) -- (axis cs:14, -5.8)
  node[right, font=\tiny, green!50!black, align=left]
    {Newton 二次收敛\\误差$^2$每步平方缩小};

% 精度参考线
\draw[gray!40, thin, dashed] (axis cs:0, -6) -- (axis cs:40, -6);
\node[right, font=\tiny, gray!50] at (axis cs:30, -6.2) {机器精度 $10^{-6}$};

% K-FAC 位置示意（介于 GD 和 Newton 之间）
\node[draw=orange!60, fill=orange!8, rounded corners=2pt,
      font=\tiny, align=center]
  at (axis cs:25, -4.5)
  {K-FAC / Shampoo\\（实用二阶方法）\\介于两者之间};

\end{axis}

% 底部说明
\node[font=\footnotesize, gray!60, align=center] at (5.0, -0.6)
  {注：L-BFGS 用 $m$ 对历史梯度近似 Hessian 逆，计算量仅 $O(mn)$；精确 Newton 需 $O(n^3)$（不可行）};

\end{tikzpicture}
\end{document}
